% ---------------------------------------------------------------------------
\section{Introduction}\label{sec:intro}
% ---------------------------------------------------------------------------
A driving policy has to do three things that pull in different directions: know a great deal about the world (what a school bus
implies, what a siren means), reason at several time scales (which road to take, which manoeuvre to start, how to move the
wheel now), and fit into a fixed power and latency envelope. Large vision--language models supply the first; the field's dominant
way of using them, autoregressive generation of a plan or trajectory \citep{hwang2024emma,renz2025simlingo,nvidia2025alpamayor1},
works against the third, because decoding one token requires streaming the model's weights once and is bound by memory bandwidth rather than arithmetic on an edge SoC.

\emph{Contrastive Language Models} (CLM) showed a different way to use a large frozen model \citep{kwok2026clm}: read the model
once (prefill), map the last hidden state and each candidate action through two small MLP heads into one shared 512-dimensional
space, and choose the action with the highest cosine score. Action embeddings are computed once and cached, the backbone is never
updated, and a three-stage curriculum trains the heads on cached features. CoVer-VLA transferred the recipe to robot manipulation with
a frozen image encoder and numeric action chunks \citep{kwok2026cover}. Both treat the action set as a flat list.

Driving has structure that a flat list wastes. Actions are \emph{hierarchical}: a route-level intent, a manoeuvre (lateral and
longitudinal, chosen separately), and a continuous trajectory. Actions are \emph{highly prior-structured}: physics, the current
speed, the last plan and the navigation command exclude most of the space before any learning. The candidate set can be very large
if it can be searched cheaply. And the input side is not fixed: not every camera or token is worth its cost at every instant, and some
decisive cues (a siren) are not visual at all.

\paragraph{Programme context.}
HiCAP is the vision--language arm of TanitAD, a programme building sub-300\,M hierarchical latent world models for driving.
The programme's newest reference planner line (anchored diffusion, REF-C) exposes the problem this paper addresses. On the parity
lineage, the best candidate in a 256-anchor fan has an ADE@2\,s of $0.1640$\,m while the candidate the planner selects has
$0.4714$\,m\evI~(REF-C-XL, 40 validation episodes; registry \S4). On the newest lineage, trained on the v7 corpus, the
hold-the-current-action control scores $0.2996$\,m and the pre-registered constant-velocity-extended control $0.2874$\,m, against $0.2975$\,m and $0.3079$\,m
for the two best trained arms\evI~(registry \S4.6--4.8; 4{,}823 windows, 141 episodes, none separated from the controls); the deficit lies
on the longitudinal axis, which carries $92.2\%$ of the gap between the planner and the hold control\evI. Strategic (route-level) quality
is reported UNAVAILABLE for two of those arms and $0.7708$ route accuracy against a $0.6742$ majority rate for the third\evI. In short, the
programme's bottleneck is neither the candidate set nor the lateral path; it is the \emph{decision}, especially the longitudinal one, and it is
under-supervised at the tactical and strategic levels. HiCAP is our attempt to build the decision maker properly on a backbone that already
knows the world.

\paragraph{Contributions.} We (1)~define a hierarchical action tree from the programme's frozen tactical/strategic token vocabulary with a factored path-by-profile trajectory vocabulary beneath it, with a retrieval-cost proposition, a beam-miss bound and a real-data pruning study (\S\ref{sec:vocab}, \S\ref{sec:theory}, \S\ref{sec:prelim});
(2)~show that a CLM-style contrastive selector is prior-biased and give a corrected, per-level objective (\S\ref{sec:train}, Corollary~\ref{cor:bias});
(3)~specify prior-encoded pruning with a conformal recall bound and a candidate-conditioned re-scoring stage that returns survivors to the backbone (\S\ref{sec:prune});
(4)~restrict imagination to what data can identify, in three tiers (soft outcome-aware targets, an analytic forecast of exogenous descriptors, a gated critic with a disagreement veto) and show in a toy that most apparent gain is target shape (\S\ref{sec:imag}, Proposition~\ref{prop:ident});
(5)~add elastic sensing (cameras, tokens, embedding width) and an audio branch tokenised into the backbone's space (\S\ref{sec:elastic});
(6)~separate a frozen reference model from a distilled sub-300\,M student with a decision-preservation guarantee (\S\ref{sec:eff}, Proposition~\ref{prop:distill});
and (7)~commit a pre-registered protocol, including negative controls and both outcomes for each hypothesis (\S\ref{sec:protocol}).

\paragraph{Scope and evidence convention.}
This is a design and pre-registration paper. \textbf{No HiCAP model has been trained.} Everything quantitative is one of: a closed-form
statement with a proof; a measurement on committed data using only inference-time quantities; a toy simulation; arithmetic; or a number
inherited from a programme document. Each number carries an evidence tag: \evM{} measured by us with an artifact in the repository,
\evP{} published and cited, \evI{} inherited from a programme document and not re-verified here, \evE{} closed-form estimate (no timing),
\evH{} hypothesis. Where two programme documents disagree, the model registry is the tie-breaker and the conflict is reported (\S\ref{sec:limits}).
